\section{Begrifssdefinitionen}
\subsection{Europäische Union}
Die Europäische Union (EU) ist ein Zusammenschluss von europäischen Staaten, die gemeinsame politische Ziele verfolgen \parencite{bpb_eu_2025}

\subsection{Amtsblatt der Europäischen Union}
Das Amtsblatt der Europäischen Union ist das Medium für die amtliche Veröffentlichung von Rechtsakten der EU, sonstigen Rechtsakten sowie amtlichen Informationen der Organe, Einrichtungen und sonstigen Stellen der EU.

Es wird montags bis freitags – und in dringenden Fällen samstags, sonntags und feiertags – in den Sprachen veröffentlicht, die zum Zeitpunkt der Veröffentlichung Amtssprachen sind (derzeit 24).

EUR-Lex enthält die elektronischen Fassungen aller Amtsblätter bis zum 30. Dezember 1952, an dem das erste Amtsblatt der Europäischen Gemeinschaft für Kohle und Stahl veröffentlicht wurde.

Seit Juli 2013 ist ausschließlich die auf EUR-Lex veröffentlichte elektronische Ausgabe des Amtsblatts verbindlich und entfaltet Rechtswirkung. \parencite{eurlex_amtsblatteu_2025}

\subsection{Europäische Richtlinie}
Eine Richtlinie ist ein Rechtsakt, in dem ein von den EU-Ländern zu erreichendes Ziel festgelegt wird. Es ist jedoch Sache der einzelnen Länder, eigene Rechtsvorschriften zur Verwirklichung dieses Ziels zu erlassen. Ein Beispiel ist die EU-Richtlinie über Einwegkunststoffe: Sie verringert die Auswirkungen bestimmter Kunststoffprodukte auf die Umwelt, indem sie beispielsweise den Gebrauch von Wegwerfplastik wie Tellern, Trinkhalmen und Getränkebechern einschränkt oder sogar verbietet. \parencite{eu_richtlinien_2025}

\subsection{Europäische Norm}
Normen und andere Veröffentlichungen zur Standardisierung sind nicht verbindliche Leitlinien mit technischen Spezifikationen für Produkte, Dienstleistungen und Verfahren – z. B. Industrieschutzhelme, Auflader für elektronische Geräte oder Dienstleistungsniveaus in öffentlichen Verkehrsmitteln. Normen werden von privaten Normungsgremien in der Regel auf Initiative von Interessenträgern entwickelt, die einen entsprechenden Bedarf ermittelt haben.

Die Einhaltung von Normen ist zwar freiwillig, stellt aber unter Beweis, dass Ihre Waren und Dienstleistungen sich durch ein bestimmtes Maß an Qualität, Sicherheit und Zuverlässigkeit auszeichnen.

Europäische Normen werden von einem der drei europäischen Normungsgremien verabschiedet:

- Europäisches Komitee für Normung CEN 

- Europäisches Komitee für elektrotechnische Normung CENELEC

- Europäisches Institut für Telekommunikationsnormen ETSI

\parencite{eu_norm_2025}

\subsection{Harmonisierte Norm}
Harmonisierte Normen sind eine besondere Kategorie europäischer Normen, deren Erarbeitung von der Europäischen Kommission bei einem europäischen Normungsgremium in Auftrag gegeben wird. Etwa 20 % aller europäischen Normen werden nach einem Normungsauftrag der Europäischen Kommission erarbeitet.

Mit der Einhaltung harmonisierter Normen stellen Sie unter Beweis, dass Ihre Produkte oder Dienstleistungen im Einklang mit den technischen Anforderungen der einschlägigen EU-Rechtsvorschriften stehen.

In den meisten Fällen ist die Anwendung harmonisierter Normen freiwillig. Als Hersteller oder Dienstleistungserbringer steht es Ihnen frei, eine andere technische Lösung zur Erfüllung der Anforderungen zu wählen.
\parencite{eu_normharmonisiert_2025}

\subsection{European Accessibility Act - EAA}
Der European Accessibility Act ist eine Richtlinie zur Verbesserung der Barrierefreiheit von Produkten und Dienstleistungen im europäischen Binnenmarkt.

Die Anforderungen an die Barrierefreiheit sind aufgrund der unterschiedlichen Regelungen in den einzelnen Mitgliedsstaaten zum Teil sehr unterschiedlich.

Der EAA soll vor allem dafür sorgen, dass in der EU gemeinsame Regeln für die Barrierefreiheit entstehen, die unter anderem zu Kostensenkungen führen, den grenzüberschreitenden Handel erleichtern und mehr Marktchancen für barrierefreie Produkte und Dienstleistungen ermöglichen. \parencite[vgl.][]{eu_eaa_2025}

\section{Gesetzliche Grundlagen und Standards}
\subsection{Barrierefreie-Informationstechnik-Verordnung (BITV 2.0)}
Die Barrierefreie-Informationstechnik-Verordnung ist am 25. Mai 2019 in Kraft getreten. Sie setzt die Vorgaben der Richtlinie 2016/2102 über die Barrierefreiheit von Websites und mobilen Anwendungen öffentlicher Stellen um. Die Vorgaben, die nicht auch schon 2018 in das aktualisierte Behindertengleichstellungsgesetz (BGG) aufgenommen wurden. 

Sie gilt für alle öffentlichen Stellen des Bundes. \parencite{bfj_bgg12_2025}

Das neue an dieser Verordnung ist, das nicht wie bisher der zu berücksichtigende Standard zur barrierefreien Gestaltung von Informationstechnik in dem Gesetz selber beschrieben wird, sondern die Verordnung auf die im Amtsblatt der Europäischen Union bekannt gemachten harmonisierten Normen verwiesen wird. 

Daraus ergibt sich ein Vorteil: Wenn die EU-Kommission in Zukunft eine andere harmonisierte Norm im Amtsblatt der EU bekannt macht, ist diese anzuwenden, ohne, dass dafür das BITV 2.0 angepasst werden muss. Dies passierte z.B. zuletzt im \parencite[vgl.][]{bbf_newversionnorm_2025}.

Wie schon erwähnt wird der anzuwendende Standard nicht direkt genannt, sondern die BITV 2.0 verweist auf die im Amtsblatt der Europäischen Union bekannt gemachten harmonisierten Normen \parencite[vgl.][]{bfj_bitv3_2025}. Diese europäische Norm ist im Moment die EN 301 549 in der Version V3.2.1 (2021-03). Im BITV 2.0 werden die Aufgaben der Überwachungsstelle festgelegt. So soll diese prüfen, ob eine Website oder App die Mindestanforderungen an Barrierefreiheit gemäß Artikel 6 der EU-Richtlinie (EU) 2016/2102 erfüllt \parencite[vgl.][]{bfj_bitv8_2025}. Die EU-Richtlinie ist vom Amtsblatt der Europäischen Union am 26. Oktober 2016 veröffentlicht worden \parencite[vgl.][]{eu_EU2016L2102_2016}.

Im Artikel 6 Absatz 2 Unterabsatz der EU-Richtlinie (EU) 2016/2102 wird dann auch die anzuwendende Norm genannt:

(2) Wurden keine Referenzen von harmonisierten Normen gemäß Absatz 1 dieses Artikels veröffentlicht, so wird bei Inhalten von mobilen Anwendungen, die die technischen Spezifikationen oder Teile davon erfüllen, davon ausgegangen, dass sie die Barrierefreiheitsanforderungen gemäß Artikel 4, die durch diese technischen Spezifikationen oder Teile davon erfasst werden, erfüllen. 
Die Kommission erlässt Durchführungsrechtsakte zur Festlegung der in Unterabsatz 1 dieses Absatzes genannten technischen Spezifikationen. Diese technischen Spezifikationen müssen die Barrierefreiheitsanforderungen gemäß Artikel 4 erfüllen und einen mit der europäischen Norm EN 301 549 V1.1.2 (2015-04) zumindest gleichwertigen Grad der Zugänglichkeit gewährleisten. \parencite{eu_EU2016L2102Norm_2016}

\subsection{Barrierefreiheitsstärkungsgesetz (BFSG)}
\subsubsection{Gesetz}
Um zu gewährleisten, dass Menschen mit Behinderungen, Einschränkungen und ältere Menschen gleichermaßen an Produkten und Dienstleistungen teilhaben können, wurde das Barrierefreiheitsstärkungsgesetz ins Leben gerufen.

Das Gesetz zur Umsetzung der Richtlinie (EU) 2019/882 des Europäischen Parlaments und des Rates über die Barrierefreiheits­anforderungen für Produkte und Dienstleistungen (BFSG) vom Bundesamt für Arbeit und Soziales (BMAS) setzt wie der Name schon sagt die Richtlinie der Europäischen Kommission (EU) 2019/882 um.

Bisher gab es für Barrierefreiheit keine einheitlichen Regelungen, an die sich die Anbieter halten mussten. In diesem Gesetz wird die EU Richtlinie zur Barrierefreiheit (European Accessibility Act, kurz: EAA) umgesetzt. Nach aktuellem Stand müssen alle, die in der Europäischen Union etwas produzieren, verkaufen oder Dienstleistungen anbieten, nicht nur ganz unterschiedliche Anforderungen für Barrierefreiheit beachten, teilweise widersprechen sich die Anforderungen sogar. 

Durch die klaren und einheitlichen Standards soll der Binnenmarkt gestärkt werden und eine größere Verfügbarkeit preisgünstiger barrierefreier Produkte und Dienstleistungen sicherstellen.

Doch warum nehmen wir auch das BITV 2.0 unter die Lupe, wenn das BFSG umgesetzt werden soll? Das BITV 2.0 ist dem BFSG sehr ähnlich. Das BITV 2.0 verpflichtet allerdings erst einmal nur die öffentlichen Stellen des Bundes die Barrierefreiheitsanforderungen umzusetzen. Mit dem BFSG werden jetzt auch private Hersteller und Dienstleister herangezogen. 

Wie schon im BITV 2.0 arbeiten die Gesetzgeber vor allem mit Verweisen. So spricht das Gesetz, wenn es um die Anforderungen geht, auch hier von harmonisierten Normen die im europäischen Amtsblatt veröffentlicht sein müssen. \parencite[vgl.][]{bmas_bfsg4_2021}. Allerdings sind zu der Richtlinie (EU) 2019/882 noch keine harmonisierten Normen vom europäischen Amtsblatt veröffentlicht worden. Dies soll aber in Zukunft passieren \parencite[vgl.][]{etsi_historyen301549_2025}.

Zur Erleichterung enthält das BFSG Konformitätsvermutungen. Werden bestimmte technische Normen, EU-harmonisierte Normen) oder DIN- oder ISO-Standards (technische Spezifikationen), erfüllt, wird die Einhaltung der Barrierefreiheitsanforderungen vermutet, wenn die Normen und Standards selbst die Barrierefreiheitsanforderungen erfüllen. Derzeit werden verschiedene neue technische Standards auf europäischer und nationaler Ebene erarbeitet

Am 01.03.2021 hat das Bundesministerium seinen ersten Referentenentwurf veröffentlicht. Am 24.03.2021 wurde er von der Regierung im Kabinett beschlossen und verabschiedet, ehe es am 22.07.2021 zum Abschluss kam.
Nach Artikel 3 tritt das Gesetz zum 28. Juni 2025 in Kraft \parencite[vgl.][]{bmas_bfsgA3_2021}.


\subsubsection{Verordnung}
Die Verordnung soll sicherstellen, dass Produkte und Dienstleistungen so gestaltet werden, dass sie für alle Menschen, insbesondere für Menschen mit Behinderungen, zugänglich und nutzbar sind. Dies umfasst sowohl technische Anforderungen als auch die Bereitstellung von Informationen und Unterstützungsdiensten. Außerdem wird auch hier noch einmal das Inkrafttreten am 28. Juni 2025 bestätigt. Die Verodnung ist am 22. Juni beschlossen worden.

\subsection{EU-Richtlinie}
Die EU-Richtlinie 2019/882 soll sicherstellen, dass Produkte und Dienstleistungen in der EU für alle Menschen, insbesondere für Menschen mit Behinderungen, leichter zugänglich sind. Ein wichtiger Teil davon ist, die teilweise sehr unterschiedlichen nationalen Anforderungen in der EU anzugleichen, um einen einheilten Standard zu garantieren. Außerdem soll die Verfügbarkeit barrierefreier Angebote erhöht und der Zugang zu wichtigen Informationen erleichtert werden. 

Die Zahl der Menschen mit Behinderung wird vorraussichtlich steigen, was die Wichtigkeit von barrierefreien Produkten und Dienstleistungen auch deutlich erhöht. Eine bessere Zugänglichkeit hilft somit Menschen mit Behinderungen, da sie ein selbstbestimmtes Leben fúhren können. 

Basierend auf der UN-Behindertenrechtskonvention, die seit 2011 für die EU verbindlich ist, gelten Menschen als behindert, wenn sie langfristige körperliche, geistige oder sensorische Einschränkungen haben, die sie in Verbindung mit Umweltfaktoren, daran hindern, gleichberechtigt am gesellschaftlichen Leben teilzunehmen. 

Daher soll die Richtlinie dafür sorgen, dass Alltagsprodukte und Dienstleistungen so gestaltet oder angepasst werden, dass sie für Menschen mit Behinderungen nutzbar sind. \parencite{eu_eur2019882begruendung_2017}

Hierbei bezieht die Richtlinie, genau wie die 2016/2102 auf der das BITV 2,0 basiert, auf 4 große Hauptpunkte \parencite[vgl.][]{eu_eur2019882Anforderungen_2017}:

- Wahrnehmbarkeit

- Bedienbarkeit

- Verständlichkeit

- Robustheit.


Genau diese vier Kategorien, sind so aus dem WCAG übernommen. Die Richtlinien und Erfolgskriterien des WCAG basieren auf den vier grundlegenden Prinzipien, die sicherstellen, dass alle Menschen, einschließlich Menschen mit Behinderungen, auf Webinhalte zugreifen und diese nutzen können. Fehlt eines dieser Prinzipien, wird die Nutzung des Webs für Menschen mit Behinderungen erheblich eingeschränkt oder unmöglich. \parencite{w3c_4principlesA11y_2024}